\appsection{Sensitivity to $(\rho, \kappa)$}
\label{app:rk_sensitivity}
%==============================================================================

This appendix reports results from re-estimating the CE-space Beta-Beta model
at three additional $(\rho, \kappa)$ configurations beyond the baseline
($\rho = 0.2$, $\kappa = 0.25$). The four configurations span the corners of
the interpretable region:

\begin{center}
\begin{tabular}{lcc}
  \toprule
  Configuration & $\rho$ & $\kappa$ \\
  \midrule
  Baseline        & 0.2 & 0.25 \\
  High $\kappa$   & 0.2 & 1.50 \\
  High $\rho$     & 0.8 & 0.25 \\
  Both high       & 0.8 & 1.50 \\
  \bottomrule
\end{tabular}
\end{center}

\noindent Low $\rho$ (0.2) allocates most direct prior variance to
within-population heterogeneity; high $\rho$ (0.8) attributes most of it to
uncertainty about the population mean. Low $\kappa$ (0.25) expresses high
confidence in the heterogeneity level; high $\kappa$ (1.50) expresses
substantial uncertainty.

All four configurations use the same direct prior (mean and variance for each
CE and the choice probability), the same data, the same Beta-Beta
distributional family, and the same diffuse NIG comparison prior used in the
main text. The only thing that changes is the
hierarchical allocation of prior variance.

\begin{center}
\begin{tabular}{lcccccc}
  \toprule
  Configuration & $\rho$ & $\kappa$ & Well & Moderate & Dominated & Dominated (\%) \\
  \midrule
  Baseline      & 0.2 & 0.25 & 79 & 14 & 18 & 16 \\
  High $\kappa$ & 0.2 & 1.50 & 83 & 14 & 14 & 13 \\
  High $\rho$   & 0.8 & 0.25 & 76 & 14 & 21 & 19 \\
  Both high     & 0.8 & 1.50 & 81 & 13 & 17 & 15 \\
  \bottomrule
\end{tabular}
\end{center}

The table makes the main pattern clear. Raising $\rho$ increases the share of
prior-dominated subjects from 16\% in the baseline to
19\% when $\rho = 0.8$ and $\kappa = 0.25$, consistent
with allocating more uncertainty to population-mean variation. Raising
$\kappa$ alone lowers the prior-dominated share to 13\%,
reflecting the weaker concentration of the hyperprior over heterogeneity.

\subsection{CDF Difference Plots}

The following figures show the CDF difference between the informative and
diffuse posteriors for each configuration, using the same format as
Figure~\ref{fig:cdf_diff} in the main text.

\begin{figure}[htbp]
  \centering
  \mycaption{CDF differences: $\rho = 0.2$, $\kappa = 1.50$}
  \label{fig:cdf_diff_rho0p2_kappa1p50}
  \includegraphics[width=0.85\textwidth]{../plots/paper/rho0p2_kappa1p50/cdf_difference_by_subject.pdf}
\end{figure}

With $\kappa$ raised to 1.50 while $\rho$ remains at 0.2, the hyperprior
over heterogeneity is substantially less concentrated. The CDF differences
are slightly attenuated relative to the baseline, consistent with the
lower prior-dominated share reported in the table above.

\begin{figure}[htbp]
  \centering
  \mycaption{CDF differences: $\rho = 0.8$, $\kappa = 0.25$}
  \label{fig:cdf_diff_rho0p8_kappa0p25}
  \includegraphics[width=0.85\textwidth]{../plots/paper/rho0p8_kappa0p25/cdf_difference_by_subject.pdf}
\end{figure}

Raising $\rho$ to 0.8 allocates the majority of direct prior variance to
uncertainty about the population mean, leaving less room for individual
heterogeneity. This produces the highest prior-dominated share across the
four configurations, with slightly larger CDF deviations for the
prior-sensitive subjects.

\begin{figure}[htbp]
  \centering
  \mycaption{CDF differences: $\rho = 0.8$, $\kappa = 1.50$}
  \label{fig:cdf_diff_rho0p8_kappa1p50}
  \includegraphics[width=0.85\textwidth]{../plots/paper/rho0p8_kappa1p50/cdf_difference_by_subject.pdf}
\end{figure}

When both $\rho$ and $\kappa$ are high, the effects partially offset.
The large $\rho$ concentrates heterogeneity while the large $\kappa$
expresses uncertainty about that concentrated specification, returning the
prior-dominated share close to the baseline level.

\subsection{Predictive CE Overlays}

The following figures compare the hyper-posterior predictive CE distributions
across configurations.

\begin{figure}[htbp]
  \centering
  \mycaption{Predictive CE overlay: $\rho = 0.2$, $\kappa = 1.50$}
  \label{fig:ce_overlay_rho0p2_kappa1p50}
  \includegraphics[width=\textwidth]{../plots/paper/rho0p2_kappa1p50/predictive_ce_overlay.pdf}
\end{figure}

\begin{figure}[htbp]
  \centering
  \mycaption{Predictive CE overlay: $\rho = 0.8$, $\kappa = 0.25$}
  \label{fig:ce_overlay_rho0p8_kappa0p25}
  \includegraphics[width=\textwidth]{../plots/paper/rho0p8_kappa0p25/predictive_ce_overlay.pdf}
\end{figure}

\begin{figure}[htbp]
  \centering
  \mycaption{Predictive CE overlay: $\rho = 0.8$, $\kappa = 1.50$}
  \label{fig:ce_overlay_rho0p8_kappa1p50}
  \includegraphics[width=\textwidth]{../plots/paper/rho0p8_kappa1p50/predictive_ce_overlay.pdf}
\end{figure}

Across all three alternatives, the overlay figures show that changing
$(\rho,\kappa)$ mostly affects the dispersion and shrinkage of the posterior
predictive distributions rather than their qualitative ordering. The same
lottery-specific pattern remains visible across configurations, with the
largest differences concentrated in the wider-range lotteries where the
informative prior has more room to move posterior mass.

\subsection{Induced Prior in Parameter Space}

The quantity-space framework specifies priors on predictive quantities and lets
the Jacobian correction induce the parameter-space prior. To visualize this
induced prior, I sample from the full hierarchical model with no data
likelihood, so that the resulting parameter draws reflect only the CE-space
specification and the Jacobian geometry.

\begin{figure}[htbp]
  \centering
  \mycaption{Hyperprior predictive distribution in parameter space}
  \label{fig:prior_predictive_params}
  \includegraphics[width=0.9\textwidth]{../plots/paper/prior_predictive_params.pdf}
\end{figure}

Figure~\ref{fig:prior_predictive_params} shows the result. The induced prior
on $r$ is centered near mild risk aversion (median 1.08) with support for both
risk-seeking and strongly risk-averse individuals. The Prelec curvature $\phi$
peaks near 1 (linear probability weighting), allowing inverse-S shapes
($\phi < 1$) that are typical of the empirical literature. The elevation
$\eta$ concentrates below 1, and the noise $\lambda$ is concentrated near zero
with a long right tail. These are sensible parameter-space priors for an RDU
model, achieved entirely through beliefs about certainty equivalents and a
choice probability.

\subsection{Structural Parameter Posteriors for Selected Subjects}

The following figures show how data updates the induced prior for individual
subjects. The most affected subject in the baseline configuration is shown in
Figure~\ref{fig:params_most}, and the least affected is shown in
Figure~\ref{fig:params_least}.

\begin{figure}[htbp]
  \centering
  \mycaption{Posterior parameter distributions: most affected subject}
  \label{fig:params_most}
  \includegraphics[width=0.9\textwidth]{../plots/paper/subject_params_most_affected.pdf}
\end{figure}

For the most affected subject, the two models recover fundamentally different
parameter regions. The informative model concentrates on high $r$ (strong risk
aversion), low $\phi$ (strong inverse-S probability weighting), and low $\eta$,
while the diffuse model favors negative $r$ (risk seeking) with different
probability weighting parameters. These are the same opposing modes visible in
the CE comparison, now shown in parameter space. The noise parameter $\lambda$
is similar across models, confirming that the disagreement is about preferences
rather than noise.

\begin{figure}[htbp]
  \centering
  \mycaption{Posterior parameter distributions: least affected subject}
  \label{fig:params_least}
  \includegraphics[width=0.9\textwidth]{../plots/paper/subject_params_least_affected.pdf}
\end{figure}

For the least affected subject, the posteriors are nearly identical across
models. The data are sufficiently informative that the prior specification has
negligible influence on the structural parameter estimates.

Taken together, these sensitivity checks reinforce the main text's
interpretation of $(\rho, \kappa)$. Changing the variance allocation and the
researcher's confidence in that allocation does move the induced prior and, for
some weakly identified subjects, can shift posterior inferences in noticeable
ways. But the broad qualitative pattern is stable across the four
configurations: the same subset of subjects remains prior-sensitive, the
lottery-specific ordering of posterior predictive CEs is preserved, and the
informative prior continues to act primarily as a regularizer where the data are
least informative. In that sense, the substantive conclusions are not tied to a
single baseline choice of $(\rho, \kappa)$, even though the magnitude of prior
influence varies in the expected direction across configurations.
